% Generated draft: requires mathematical and source review.
\documentclass[11pt]{article}
\usepackage[margin=25mm]{geometry}
\usepackage{fontspec}
\setmainfont{Latin Modern Roman}
\usepackage{amsmath,amssymb,mathtools,unicode-math}
\setmathfont{Latin Modern Math}
\usepackage{enumitem,needspace,xurl,hyperref}
\setlength{\parindent}{0pt}
\setlength{\parskip}{5pt}
\setlength{\emergencystretch}{3em}
\newcommand{\N}{\mathbb{N}}
\newcommand{\Z}{\mathbb{Z}}
\newcommand{\Q}{\mathbb{Q}}
\newcommand{\R}{\mathbb{R}}
\newcommand{\C}{\mathbb{C}}
\newcommand{\F}{\mathbb{F}}
\newcommand{\A}{\mathbb{A}}
\newcommand{\PP}{\mathbb{P}}
\newcommand{\E}{\mathbb{E}}
\newcommand{\eps}{\varepsilon}
\newcommand{\dd}{\,\mathrm{d}}
\begin{document}
\section*{UnsolvedMath ID: 700198}
\textbf{Source problem code:} catalogue entry
\subsection{Definitions and mathematical statement}
A semigroup is a nonempty collection equipped with an associative binary operation. For a semigroup \(S\), a nontrivial semigroup identity is an equation \(u=v\), where \(u\) and \(v\) are distinct nonempty words in finitely many variables \(x_1,\ldots,x_n\), such that
\[
\forall (s_1,\ldots,s_n)\in S^n,\qquad u(s_1,\ldots,s_n)=v(s_1,\ldots,s_n).
\]
Here word evaluation uses the semigroup operation, and no identity element is assumed.

The supplied question of B. I. Plotkin asks whether the semigroup generated by all indecomposable radicals satisfies any identity. Conditional on specifying the intended radicals, their indecomposability condition, and their associative multiplication, let \(S\) denote the semigroup of their finite nonempty products. Interpreting ``any identity'' as a nontrivial semigroup identity, the question is whether
\[
\exists n\geq 1\;\exists u\neq v\text{ nonempty words in }x_1,\ldots,x_n\quad
\forall (s_1,\ldots,s_n)\in S^n,\quad u(s_1,\ldots,s_n)=v(s_1,\ldots,s_n).
\]
The supplied record does not define ``radical'', ``indecomposable'', or the multiplication of radicals. Consequently, it does not identify \(S\) sufficiently to give an unconditional mathematical formulation; the displayed statement is a conditional formulation of the requested identity question.
\subsection{Short English statement}
B. I. Plotkin asks whether all finite nonempty products of indecomposable radicals obey a common nontrivial semigroup identity. The intended radicals, indecomposability condition, and multiplication must be specified before this question is fully determined.
\subsection{Sources}
\begin{itemize}
\item Supplied local UnsolvedMath registry record: permanent ID \(700198\), exact source code \texttt{catalogue entry}, title \emph{Kourovka Notebook Problem 3.49}; upstream revision \texttt{d371ee917ac9647fbb06cd6fc529deec6f3dda82}.
\item The record attributes Problem 3.49 to B. I. Plotkin, proposed in 1969, and identifies its source as \emph{The Kourovka Notebook (Unsolved Problems in Group Theory)}, 3rd issue (1969), with the 21st issue edition (2026), \texttt{arXiv:1401.0300v48}. The supplied collection description names E. I. Khukhro and V. D. Mazurov as editors.
\item The supplied registry labels the problem \emph{open}. This is a supplied status label, not an independently verified mathematical or literature assessment. The cited source contents were not supplied beyond the registry quotation.
\end{itemize}
\end{document}
